% ---------------------------------------------------------------------------
% Proyecto I - Robot Aspiradora Autonomo (Yocto)
% CE-1113 Sistemas Empotrados - TEC
%
% Diagramas de hardware. El detalle (calculos, listas de materiales y
% procedimientos de verificacion) vive en docs/hardware-*.md; aqui solo
% estan los diagramas.
%
% Compilar:  pdflatex hardware-diagramas.tex   (dos veces)
% ---------------------------------------------------------------------------
\documentclass[11pt,a4paper]{article}

\usepackage[T1]{fontenc}
\usepackage[utf8]{inputenc}
\usepackage{lmodern}
\usepackage[spanish,es-noshorthands,es-noquoting]{babel}
\usepackage[margin=2.4cm]{geometry}
\usepackage{tikz}
\usepackage[siunitx,RPvoltages]{circuitikz}
\usepackage{amsmath}
\usepackage{amssymb}
\usepackage{booktabs}
\usepackage{caption}

\usetikzlibrary{positioning,arrows.meta,fit,backgrounds,calc,shapes.geometric}

\captionsetup{font=small,labelfont=bf,width=.92\linewidth}
\setlength{\parskip}{0.4em}
\setlength{\parindent}{0pt}

% --- paleta -----------------------------------------------------------------
\definecolor{cLog}{RGB}{31,102,168}
\definecolor{cPot}{RGB}{193,84,26}
\definecolor{cLed}{RGB}{34,120,64}
\definecolor{cGris}{RGB}{120,120,120}
\definecolor{cFondo}{RGB}{245,247,250}

% --- estilos comunes --------------------------------------------------------
\tikzset{
  blk/.style   ={draw,rounded corners=2pt,align=center,inner sep=4pt,
                 minimum height=8mm,font=\small},
  blkl/.style  ={blk,draw=cLog,fill=cLog!8},
  blkp/.style  ={blk,draw=cPot,fill=cPot!8},
  blkn/.style  ={blk,draw=cGris,fill=black!3},
  rot/.style   ={font=\scriptsize\bfseries},
  ln/.style    ={-{Latex[length=2mm]},thick},
}

% --- caja de nota -----------------------------------------------------------
\newsavebox{\notabox}
\newenvironment{nota}
  {\par\medskip\begin{lrbox}{\notabox}%
   \begin{minipage}{\dimexpr0.95\linewidth-14pt\relax}\small}
  {\end{minipage}\end{lrbox}%
   \begin{center}\setlength{\fboxsep}{7pt}\colorbox{cFondo}{\usebox{\notabox}}\end{center}\medskip}

% ---------------------------------------------------------------------------
\title{\vspace{-1.4cm}\textbf{Diagramas de hardware}\\[2pt]
       \large Proyecto I --- Robot Aspiradora Autónomo}
\author{\small CE-1113 Sistemas Empotrados --- Tecnológico de Costa Rica}
\date{\small\today}

\begin{document}
\maketitle
\vspace{-0.8cm}

{\small Los cálculos, las listas de materiales y los procedimientos de
verificación están en \texttt{docs/hardware-*.md}. Este documento reúne
únicamente los diagramas.}

\vspace{0.2cm}
\hrule
\vspace{0.4cm}

% ===========================================================================
\section{Arquitectura general}

Un pack 2S alimenta dos rieles regulados por separado. Las seis señales de
control del puente H cruzan hacia el dominio de potencia a través de
optoacopladores.

\begin{figure}[h]
\centering
\begin{tikzpicture}[font=\small]

  % --- fuente ---
  \node[blkn] (bat) at (0,5.4) {Pack 2S 18650\\[-1pt]\scriptsize 7{,}4\,V nom. (6{,}0--8{,}4\,V)};
  \node[blkn,right=8mm of bat] (bms) {BMS 2S\\[-1pt]\scriptsize $\ge$\,10\,A, con balanceo};
  \draw[ln] (bat) -- (bms);

  % --- nodo de reparto ---
  \draw[thick] (bms.south) -- (bms.south |- 0,3.3);
  \draw[thick] (-3.1,3.3) -- (5.1,3.3);
  \coordinate (spL) at (-3.1,3.3);
  \coordinate (spR) at (5.1,3.3);

  % --- reguladores ---
  \node[blkl] (buck) at (-3.1,2.35) {Buck 5\,V\\[-1pt]\scriptsize MP2307, 3\,A};
  \node[blkp] (dir)  at (5.1,2.35)  {Directo\\[-1pt]\scriptsize 7{,}4\,V};
  \draw[ln] (spL) -- (buck);
  \draw[ln] (spR) -- (dir);

  % --- etiquetas de riel ---
  \node[rot,text=cLog,below=2mm of buck] (rl) {RIEL LÓGICO\ \ 5\,V / GND\_LOG};
  \node[rot,text=cPot,below=2mm of dir]  (rp) {RIEL POTENCIA\ \ 7{,}4\,V / GND\_POT};

  % --- consumidores ---
  \node[blkl,below=3mm of rl,text width=3.5cm] (pi)
    {\textbf{Raspberry Pi 4 B}\\[-1pt]\scriptsize GPIO (BCM) $\cdot$ jack 3{,}5\,mm};
  \node[blkl,below=2.5mm of pi,text width=3.5cm] (per)
    {\scriptsize 3 $\times$ HC-SR04\\[-1pt]\scriptsize 4 $\times$ LED indicador\\[-1pt]\scriptsize LM386 + altavoz 8\,$\Omega$};
  \draw[thick,cLog] (pi) -- (per);

  \node[blkp,below=3mm of rp,text width=3.3cm] (l298)
    {\textbf{Driver L298N}\\[-1pt]\scriptsize doble puente H};
  \node[blkp,below=2.5mm of l298,text width=3.3cm] (mot)
    {\scriptsize Motor izq. \quad Motor der.\\[-1pt]\scriptsize DC con reductora};
  \draw[ln,cPot] (l298) -- (mot);

  % --- barrera ---
  \node[blk,draw=cLed,fill=cLed!8,text=cLed] (opto) at (1.0,-0.15)
    {6 $\times$ PC817\\[-1pt]\scriptsize ENA, IN1, IN2,\\[-2pt]\scriptsize ENB, IN3, IN4};
  \draw[ln] (pi.east)   -- (opto.west);
  \draw[ln] (opto.east) -- (l298.west);

  \draw[dashed,thick,cLed] (1.0,1.5) -- (1.0,-2.0);
  \node[rot,text=cLed,fill=white,inner sep=2pt] at (1.0,1.05) {AISLAMIENTO ÓPTICO};

  % --- retorno comun ---
  \draw[dotted,thick,cGris] (buck.west) -- ++(-0.75,0) |- (-3.85,-2.55) -- (5.85,-2.55) |- (dir.east);
  \node[font=\scriptsize\itshape,text=cGris] at (1.0,-2.8)
    {retorno común del pack --- ver nota};

\end{tikzpicture}
\caption{Arquitectura de hardware: fuente única, dos rieles regulados por
separado y barrera óptica en las seis señales de control.}
\end{figure}

\begin{nota}
\textbf{Sobre la separación de tierras.} El MP2307 es un reductor
\emph{no aislado}: su tierra de salida es la misma de entrada. Con los dos
rieles colgando del mismo pack, \texttt{GND\_LOG} y \texttt{GND\_POT} quedan
unidos en el negativo de la batería (trazo punteado), y la prueba de
«circuito abierto» del Paso 4 de \texttt{docs/hardware-alimentacion.md} no
puede pasar tal cual. Para separación real hacen falta \emph{dos fuentes}:
un convertidor DC-DC aislado (p.\,ej.\ B0505S-1W) para el riel lógico, o un
pack aparte para los motores. Los optoacopladores siguen valiendo la pena
—cortan el lazo de ruido por la señal—, pero conviene decidir esto antes de
armar la placa.
\end{nota}

\clearpage
% ===========================================================================
\section{Aislamiento de las señales de control}

Un canal, replicado seis veces (\texttt{ENA}, \texttt{IN1}, \texttt{IN2},
\texttt{ENB}, \texttt{IN3}, \texttt{IN4}).

\begin{figure}[h]
\centering
\begin{circuitikz}[scale=1.0,transform shape,font=\small]

  % --- rotulos de dominio ---
  \node[font=\scriptsize\bfseries,text=cLog,anchor=east] at (3.0,4.6)
        {DOMINIO LÓGICO (3{,}3\,V)};
  \node[font=\scriptsize\bfseries,text=cPot,anchor=west] at (3.3,4.6)
        {DOMINIO DE POTENCIA (7{,}4\,V)};

  % --- lado logico: GPIO -> R -> LED del opto ---
  \draw (-0.5,1.6) node[left]{GPIO} to[R=$390\,\Omega$] (2.2,1.6);
  \draw (2.2,1.6) to[full led] (2.2,0.2);
  \draw (2.2,0.2) -- (2.2,-0.15) node[ground]{};
  \node[font=\scriptsize,anchor=east] at (2.95,-0.85) {GND\_LOG};

  % --- encapsulado del opto ---
  \draw[dashed,rounded corners=2pt,cLed,thick] (1.7,-0.05) rectangle (4.9,2.2);
  \node[font=\scriptsize\bfseries,text=cLed] at (3.15,2.45) {PC817};

  % --- lado potencia: fototransistor + pull-up ---
  \draw (4.3,0.9) node[npn,anchor=B,scale=0.9] (Q) {};
  \draw (Q.C) -- ++(0,1.0) coordinate (ct);
  \draw (ct) to[R=$4{,}7\,\mathrm{k}\Omega$] ($(ct)+(0,1.3)$)
        node[above,font=\scriptsize]{$+5\,\mathrm{V}\_\mathrm{POT}$};
  \draw (Q.E) -- (Q.E |- 0,-0.15) node[ground]{};
  \node[font=\scriptsize,anchor=west] at (3.35,-0.85) {GND\_POT};
  \draw (Q.C) -- ++(2.0,0) node[right,align=left,font=\scriptsize]
        {entrada del\\ L298N};

  % --- barrera ---
  \draw[dashed,thick,cLed] (3.15,-1.15) -- (3.15,4.3);

\end{circuitikz}
\caption{Canal de aislamiento con PC817 (1 de 6). Emisor común: el
fototransistor conduce cuando el LED enciende, de modo que \textbf{la señal
sale invertida}. Se compensa en \texttt{lib/lib\_motors.c}.}
\end{figure}

\begin{center}
\small
\begin{tabular}{lccc}
\toprule
\textbf{GPIO} & \textbf{LED del opto} & \textbf{Fototransistor} & \textbf{Entrada L298N}\\
\midrule
Alto (3{,}3\,V) & encendido & conduce & \textbf{Bajo}\\
Bajo (0\,V)     & apagado   & en corte & \textbf{Alto} (por el \emph{pull-up})\\
\bottomrule
\end{tabular}
\end{center}

\begin{nota}
\textbf{Consecuencia crítica.} Sin compensar la inversión,
\texttt{motores\_detener()} escribe ciclo 0 y entrega \textbf{velocidad
máxima}. La corrección va en software (\texttt{PWM\_DUTY(d) = 255 - d} y
\texttt{GPIO\_NIVEL(v) = !v}) y debe estar verificada antes de la primera
prueba con motores montados.
\end{nota}

\clearpage
% ===========================================================================
\section{Acondicionamiento de señal}

\begin{figure}[h]
\centering
\begin{minipage}[t]{0.46\linewidth}
\centering
\begin{circuitikz}[scale=0.92,transform shape,font=\small]
  \draw (0,3.4) node[above,font=\scriptsize,align=center]{HC-SR04 \texttt{ECHO}\\ (5\,V)}
        -- (0,3.0)
        to[R=$1\,\mathrm{k}\Omega$] (0,1.6);
  \draw (0,1.6) -- (1.6,1.6) node[right,font=\scriptsize,align=left]
        {GPIO\\ (3{,}3\,V máx.)};
  \draw (0,1.6) to[R=$2\,\mathrm{k}\Omega$] (0,0.1);
  \draw (0,0.1) node[ground]{};
  \node[font=\scriptsize,anchor=west] at (0.35,-0.5) {GND\_LOG};
\end{circuitikz}
\\[4pt]
{\footnotesize $V_{out}=5\,\mathrm{V}\cdot\dfrac{2\,\mathrm{k}}{1\,\mathrm{k}+2\,\mathrm{k}}=3{,}33\,\mathrm{V}$}
\\[2pt]
{\scriptsize \textbf{(a)} Divisor en \texttt{ECHO} --- 3 unidades, una por sensor.}
\end{minipage}
\hfill
\begin{minipage}[t]{0.46\linewidth}
\centering
\begin{circuitikz}[scale=0.92,transform shape,font=\small]
  \draw (0,1.6) node[left,font=\scriptsize]{GPIO}
        to[R=$R$] (2.4,1.6)
        to[full led] (2.4,0.3);
  \draw (2.4,0.3) -- (2.4,0.1) node[ground]{};
  \node[font=\scriptsize,anchor=west] at (2.75,-0.5) {GND\_LOG};
\end{circuitikz}
\\[4pt]
{\footnotesize $R=(3{,}3\,\mathrm{V}-V_f)/I_f$, \ con $I_f{=}5\,\mathrm{mA}$}
\\[2pt]
{\scriptsize \textbf{(b)} LED indicador --- 4 unidades.}
\end{minipage}
\caption{Acondicionamiento entre el GPIO y el mundo exterior. El divisor
\textbf{no es opcional}: los GPIO toleran 3{,}3\,V como máximo absoluto y no
tienen protección de entrada.}
\end{figure}

\begin{center}
\small
\begin{tabular}{lcccc}
\toprule
\textbf{LED} & \textbf{Color} & \textbf{$V_f$} & \textbf{$R$ comercial} & \textbf{GPIO (BCM)}\\
\midrule
Encendido & verde    & 2{,}2\,V & $220\,\Omega$ & 16\\
Autónomo  & azul     & 3{,}0\,V & $68\,\Omega$  & 20\\
Manual    & amarillo & 2{,}1\,V & $270\,\Omega$ & 21\\
Obstáculo & rojo     & 2{,}0\,V & $270\,\Omega$ & 26\\
\bottomrule
\end{tabular}
\end{center}

% ===========================================================================
\section{Salida de audio}

Jack analógico de 3{,}5\,mm más un LM386. No consume GPIO, que están al
límite del presupuesto de corriente.

\begin{figure}[h]
\centering
\begin{circuitikz}[scale=0.9,transform shape,font=\small]

  % --- encapsulado ---
  \draw[rounded corners=2pt,fill=black!4,draw=black!55] (3.0,0.5) rectangle (5.7,3.1);
  \node[font=\small\bfseries] at (4.35,2.72) {LM386};
  \node[font=\scriptsize,anchor=west] at (3.08,2.15) {2\ \ $+$in};
  \node[font=\scriptsize,anchor=west] at (3.08,1.45) {3\ \ $-$in};
  \node[font=\scriptsize,anchor=west] at (3.08,0.85) {4\ \ GND};
  \node[font=\scriptsize,anchor=east] at (5.62,2.15) {6\ \ $V_s$};
  \node[font=\scriptsize,anchor=east] at (5.62,1.45) {5\ \ out};
  \node[font=\scriptsize,anchor=east] at (5.62,0.85) {1--8\ gain};

  % --- entrada y volumen ---
  \draw (-0.7,2.15) node[left,align=right,font=\scriptsize]{jack 3{,}5\,mm\\(canal L)}
        to[C=$10\,\mu\mathrm{F}$] (1.5,2.15) -- (3.0,2.15);
  \draw (1.5,2.15) to[pR,l_=$10\,\mathrm{k}\Omega$] (1.5,0.35) node[ground]{};
  \node[font=\tiny,anchor=west] at (1.72,0.62) {volumen};
  \draw (3.0,1.45) -- (2.45,1.45) -- (2.45,0.75) node[ground]{};
  \draw (4.35,0.5) -- (4.35,0.05) node[ground]{};
  \node[font=\scriptsize] at (4.35,-0.8) {GND\_LOG};

  % --- alimentacion y desacople ---
  \draw (5.7,2.15) -- (6.4,2.15) -- (6.4,3.9)
        node[above,font=\scriptsize]{$+5\,\mathrm{V}\_\mathrm{LOG}$};
  \draw (6.4,2.85) to[C=$100\,\mu\mathrm{F}$,*-] (8.5,2.85) -- (8.5,2.45) node[ground]{};

  % --- salida ---
  \draw (5.7,1.45) -- (6.3,1.45) to[C=$220\,\mu\mathrm{F}$] (7.9,1.45);
  \draw (7.9,1.45) to[loudspeaker] (7.9,0.15);
  \node[font=\scriptsize,anchor=east] at (7.45,0.8) {8\,$\Omega$};
  \draw (7.9,0.15) -- (7.9,-0.1) node[ground]{};

  % --- red de Zobel ---
  \draw (7.9,1.45) -- (9.5,1.45) to[R=$10\,\Omega$] (9.5,0.7)
        to[C=$47\,\mathrm{nF}$] (9.5,-0.1) node[ground]{};
  \node[font=\scriptsize,anchor=west,align=left] at (9.8,2.0) {red de\\Zobel};

\end{circuitikz}
\caption{Amplificador de audio. Se alimenta del \textbf{riel lógico}, no del
de potencia, para que el ruido de conmutación de los motores no entre al
audio. El potenciómetro fija el nivel máximo; el volumen del enunciado se
controla por software con \texttt{amixer}.}
\end{figure}

\clearpage
% ===========================================================================
\section{Mapa del conector de 40 pines}

Numeración BCM entre paréntesis. Los pines sin etiqueta quedan libres.

\begin{figure}[h]
\centering
\begin{tikzpicture}[font=\scriptsize,y=-0.435cm]

  \def\dx{1.05}

  % --- cuerpo del conector ---
  \draw[rounded corners=3pt,fill=black!4,draw=black!35]
        (-\dx-0.42,-0.55) rectangle (\dx+0.42,19.55);

  % --- pines libres / alimentacion, columna impar (izquierda) ---
  \foreach \p/\lbl/\c in {%
      1/3V3/cGris, 9/GND/cGris, 11/{TRIG frontal (17)}/cLog,
      13/{ECHO frontal (27)}/cLog, 15/{TRIG lat. izq. (22)}/cLog,
      17/3V3/cGris, 19/{ECHO lat. izq. (10)}/cLog,
      21/{TRIG lat. der. (9)}/cLog, 23/{ECHO lat. der. (11)}/cLog,
      25/GND/cGris, 29/{IN1 motor izq. (5)}/cPot,
      31/{IN2 motor izq. (6)}/cPot, 33/{ENB PWM der. (13)}/cPot,
      37/{LED obstáculo (26)}/cLed, 39/GND/cGris}
  {
    \pgfmathtruncatemacro{\r}{(\p-1)/2}
    \fill[\c] (-\dx,\r) circle (2.1pt);
    \node[anchor=east,text=\c] at (-\dx-0.55,\r) {\lbl};
  }

  % --- columna par (derecha) ---
  \foreach \p/\lbl/\c in {%
      2/5V/cGris, 4/5V/cGris, 6/GND/cGris, 14/GND/cGris,
      16/{IN3 motor der. (23)}/cPot, 18/{IN4 motor der. (24)}/cPot,
      20/GND/cGris, 30/GND/cGris, 32/{ENA PWM izq. (12)}/cPot,
      34/GND/cGris, 36/{LED encendido (16)}/cLed,
      38/{LED autónomo (20)}/cLed, 40/{LED manual (21)}/cLed}
  {
    \pgfmathtruncatemacro{\r}{(\p-2)/2}
    \fill[\c] (\dx,\r) circle (2.1pt);
    \node[anchor=west,text=\c] at (\dx+0.55,\r) {\lbl};
  }

  % --- pines libres ---
  \foreach \p in {3,5,7,27,35}
    { \pgfmathtruncatemacro{\r}{(\p-1)/2}
      \draw[black!30] (-\dx,\r) circle (2.1pt); }
  \foreach \p in {8,10,12,22,24,26,28}
    { \pgfmathtruncatemacro{\r}{(\p-2)/2}
      \draw[black!30] (\dx,\r) circle (2.1pt); }

  % --- numeros de pin ---
  \foreach \r in {0,...,19} {
    \pgfmathtruncatemacro{\a}{2*\r+1}
    \pgfmathtruncatemacro{\b}{2*\r+2}
    \node[text=black!45,font=\tiny,anchor=east] at (-0.18,\r) {\a};
    \node[text=black!45,font=\tiny,anchor=west] at (0.18,\r) {\b};
  }

  % --- leyenda ---
  \node[anchor=west,text=cLog,font=\scriptsize\bfseries]  at (-5.6,20.9) {$\bullet$\ \ sensores};
  \node[anchor=west,text=cPot,font=\scriptsize\bfseries]  at (-3.1,20.9) {$\bullet$\ \ motores (vía opto)};
  \node[anchor=west,text=cLed,font=\scriptsize\bfseries]  at (0.6,20.9) {$\bullet$\ \ LEDs};
  \node[anchor=west,text=cGris,font=\scriptsize\bfseries] at (2.9,20.9) {$\bullet$\ \ alimentación};

\end{tikzpicture}
\caption{Conector de 40 pines. GPIO 12 y 13 son los canales PWM0/PWM1 del
SoC: permiten migrar a PWM por hardware sin recablear. GPIO 14/15 (UART)
quedan libres para la consola serie de depuración.}
\end{figure}

\clearpage
% ===========================================================================
\section{Modelo físico}

Chasis circular de dos niveles. La forma circular permite girar sobre el
propio eje sin barrer área adicional: cualquier rincón del que el robot pueda
entrar es un rincón del que puede salir.

\begin{figure}[h]
\centering
\begin{minipage}[b]{0.47\linewidth}
\centering
\begin{tikzpicture}[font=\scriptsize,scale=1.0]

  \draw[fill=black!3,draw=black!45,thick] (0,0) circle (2.5);

  % --- sensores ---
  \foreach \a/\lbl in {90/frontal, 135/{lat. izq.}, 45/{lat. der.}} {
    \draw[fill=cLog!22,draw=cLog,thick]
      ([shift={(\a:2.32)}]0,0) ++(\a:0) circle (0.2);
    \draw[cLog,thick,-{Latex[length=1.6mm]}]
      (\a:2.55) -- (\a:3.25);
  }
  \node[text=cLog,anchor=south] at (90:3.35) {frontal 0$^\circ$};
  \node[text=cLog,anchor=east]  at (135:3.35) {lat. 45$^\circ$};
  \node[text=cLog,anchor=west]  at (45:3.35)  {lat. 45$^\circ$};

  % --- eje de traccion, ligeramente adelante del centro ---
  \draw[cPot,thick,dashed] (-2.5,0.45) -- (2.5,0.45);
  \draw[fill=cPot!25,draw=cPot,thick] (-2.62,0.05) rectangle (-1.85,0.85);
  \draw[fill=cPot!25,draw=cPot,thick] ( 1.85,0.05) rectangle ( 2.62,0.85);
  \node[text=cPot,anchor=east,align=right] at (-2.75,0.45) {motor\\ izq.};
  \node[text=cPot,anchor=west,align=left] at (2.75,0.45) {motor\\ der.};

  % --- centro geometrico ---
  \fill[black!55] (0,0) circle (1.4pt);
  \node[anchor=north west,text=black!55,font=\tiny] at (0.08,-0.05) {centro};

  % --- rueda loca ---
  \draw[fill=black!12,draw=black!50] (0,-1.75) circle (0.28);
  \node[anchor=south,font=\tiny] at (0,-1.42) {rueda loca};

  % --- diametro ---
  \draw[{Latex[length=1.6mm]}-{Latex[length=1.6mm]},black!55]
       (-2.5,-2.95) -- (2.5,-2.95);
  \node[anchor=north,text=black!55] at (0,-3.0) {$\varnothing$ 22--25\,cm};

\end{tikzpicture}
\\[3pt]
{\scriptsize \textbf{(a)} Vista superior. El eje va algo adelante del centro
para cargar peso sobre las ruedas motrices.}
\end{minipage}
\hfill
\begin{minipage}[b]{0.47\linewidth}
\centering
\begin{tikzpicture}[font=\scriptsize,scale=0.86,transform shape]

  % --- piso ---
  \draw[thick,black!55] (-3.0,0) -- (3.2,0);
  \foreach \x in {-2.85,-2.45,...,3.05}
    \draw[black!35] (\x,0) -- ++(-0.16,-0.16);

  % --- placas ---
  \draw[fill=black!8,draw=black!55,thick] (-2.5,0.95) rectangle (2.5,1.10);
  \draw[fill=black!8,draw=black!55,thick] (-2.5,2.90) rectangle (2.5,3.05);

  % --- separadores M3 x 40 mm ---
  \foreach \x in {-2.25,2.25}
    \draw[black!50,thick] (\x,1.10) -- (\x,2.90);

  % --- contenido de cada nivel ---
  \node[font=\tiny,anchor=south,align=center] at (0,3.18)
    {\textbf{nivel superior:} Pi 4, LEDs,\\ altavoz, interruptor};
  \node[font=\tiny,align=center] at (0,2.00)
    {\textbf{nivel inferior:} pack 2S + BMS,\\ buck 5\,V, optos + L298N, LM386};

  % --- rueda motriz ---
  \draw[fill=cPot!20,draw=cPot,thick] (-1.35,0.48) circle (0.48);
  \node[text=cPot,anchor=east,font=\tiny] at (-2.10,0.48) {$\varnothing$\,6{,}5\,cm};

  % --- sensor en el borde ---
  \draw[fill=cLog!22,draw=cLog,thick] (2.5,1.25) rectangle (2.78,1.70);
  \draw[cLog,-{Latex[length=1.5mm]}] (2.88,1.48) -- (3.30,1.48);

  % --- cotas ---
  \draw[{Latex[length=1.5mm]}-{Latex[length=1.5mm]},black!55] (-3.30,0) -- (-3.30,3.05);
  \node[anchor=east,text=black!55,font=\tiny,align=right] at (-3.37,1.90) {10--12\\cm};

  \draw[{Latex[length=1.5mm]}-{Latex[length=1.5mm]},cLog] (3.45,0) -- (3.45,1.48);
  \node[anchor=west,text=cLog,font=\tiny,align=left] at (3.52,0.74) {5--8\\cm};

  \draw[{Latex[length=1.5mm]}-{Latex[length=1.5mm]},black!55] (0.35,0) -- (0.35,0.95);
  \node[anchor=west,text=black!55,font=\tiny] at (0.45,0.48) {1{,}5--2\,cm};

\end{tikzpicture}
\\[3pt]
{\scriptsize \textbf{(b)} Vista lateral. La batería va al centro y lo más
bajo posible: baja el centro de gravedad y evita el cabeceo.}
\end{minipage}
\caption{Modelo físico. Masa objetivo bajo 1{,}5\,kg; por encima, los motores
DC pequeños pierden par para arrancar.}
\end{figure}

\end{document}
